\section{Evaluasi}
\label{sec:evaluasi-bab17}

Evaluasi pada bab ini berbentuk \textbf{ujian komprehensif}. Soal
menuntut penggabungan konsep dari beberapa bab sekaligus.

\subsection{Soal Komprehensif}

\begin{enumerate}
    \item \textbf{Analisis terpadu.} Diberikan potongan program
          berikut.

\begin{lstlisting}[language=MASM, numbers=none]
.MODEL SMALL
.STACK 100h
.DATA
    data1 DW 0A000h
    data2 DW 0004h
.CODE
START:
    MOV AX, @DATA
    MOV DS, AX
    MOV AX, data1
    MOV BX, data2
    ADD AX, BX
    MOV AH, 4Ch
    INT 21h
END START
\end{lstlisting}

          Jawab pertanyaan berikut.
          \begin{enumerate}
              \item Berapa nilai \reg{AX} setelah \asm{ADD AX, BX}?
              \item Nyatakan nilai tersebut dalam heksadesimal dan
                    tafsiran bertanda.
              \item Tentukan \flag{CF}, \flag{ZF}, \flag{SF}, dan
                    \flag{OF} beserta alasannya.
              \item Bila \code{data2} diubah menjadi \hex{C000},
                    flag mana yang berubah dan mengapa?
          \end{enumerate}

    \item \textbf{Perancangan prosedur.} Rancang prosedur
          \code{RataRata} yang menerima tiga bilangan 16-bit melalui
          stack dan mengembalikan rata-ratanya pada \reg{AX}.
          Tuliskan gambar isi stack, kode prosedur dengan
          \eng{prologue}/\eng{epilogue}, dan cara pemanggilannya.

    \item \textbf{Perbaikan bug.} Program bermaksud menjumlahkan 10
          elemen array, tetapi hasilnya selalu kurang satu elemen.
          Tuliskan langkah sistematis menemukan penyebabnya dan
          perintah \eng{debugger} yang dipakai.

    \item \textbf{Array, string, dan pencetakan.} Diketahui array
          word dan string. Tuliskan rancangan program yang
          menjumlahkan array, menyalin string, lalu mencetak hasil
          penjumlahan dalam desimal. Sebutkan prosedur dan register
          yang dipakai.

    \item \textbf{Interupsi.} Jelaskan urutan penyimpanan dan
          pemulihan \reg{FLAGS}, \reg{CS}, dan \reg{IP} ketika
          \code{INT 21h} dan \code{IRET} dijalankan. Sertakan diagram
          stack.

    \item \textbf{Perancangan.} Anda diminta membuat program yang
          menampilkan isi array dalam bentuk heksadesimal. Tuliskan
          daftar prosedur beserta masukan dan keluarannya sebelum
          menulis kode.
\end{enumerate}

\subsection{Kisi-kisi Ujian Komprehensif}

\begin{table}[H]
    \centering
    \caption{Kisi-kisi ujian komprehensif
    Bab~\ref{chap:evaluasi-akhir}.}
    \label{tab:kisi-bab17}
    \small
    \begin{tabular}{@{}p{4.0cm}p{5.4cm}c@{}}
        \toprule
        \textbf{Aspek yang diukur} & \textbf{Soal} & \textbf{Bobot} \\
        \midrule
        Integrasi aritmatika, representasi, dan \eng{flag} & Soal 1 &
        20\% \\
        \addlinespace
        Perancangan dan penulisan prosedur & Soal 2, 6 & 25\% \\
        \addlinespace
        Metodologi \eng{debugging} & Soal 3 & 15\% \\
        \addlinespace
        Array, string, dan pencetakan bilangan & Soal 4 & 20\% \\
        \addlinespace
        Mekanisme interupsi & Soal 5 & 20\% \\
        \bottomrule
    \end{tabular}
\end{table}

\begin{checklistbox}
\vspace{1.1em}
\textbf{Checklist Kesiapan Menghadapi UAS}

\begin{itemize}
    \item[\checkbox] Saya dapat menulis program aritmatika dan
          manipulasi bit tanpa melihat contoh.
    \item[\checkbox] Saya dapat menerjemahkan pseudocode bercabang
          dan berulang menjadi assembly.
    \item[\checkbox] Saya dapat menggambar isi stack dan menghitung
          parameter melalui \reg{BP}.
    \item[\checkbox] Saya dapat menelusuri program dengan
          \eng{debugger} dan mengisi tabel \eng{trace}.
    \item[\checkbox] Saya dapat mengolah array/string dan mencetak
          bilangan.
    \item[\checkbox] Saya dapat menjelaskan mekanisme interupsi
          beserta perubahannya pada stack.
\end{itemize}
\end{checklistbox}
